\begin{textsample}{NEG Dim 3 – 2009\_c6 – Score -3.00 – 2009\_c6/enade\_2009\_com\_soc\_edit\_bd87747c88.txt}  \label{ex:f3_neg_049}
Capítulo 3 Análise Técnica da Prova Este capítulo tem por objetivo apresentar o desempenho dos estudantes de COMUNICAÇÃO SOCIAL – EDITORAÇÃO no ENADE/2009.
Para isso, foram calculadas as estatísticas básicas da prova como um todo, bem como as estatísticas das questões relacionadas à Formação Geral e ao Componente Específico.
Nas Tabelas, são apresentadas as seguintes estatísticas : total da população e de presentes, média, erro-padrão da média, desvio-padrão, mediana, nota máxima, nota mínima e coeficiente de assimetria ( ver definição em 1.4.3 ).
Nas estatísticas das questões discursivas, foram detalhadas também as notas nulas : em branco, por protesto ou por insuficiência.
As estatísticas apresentadas neste capítulo contemplam, separadamente, os ingressantes, os concluintes e o total de estudantes, e foram calculadas tendo -se em vista as seguintes agregações : ( a ) Região Geográfica e País e ( b ) a Categoria Administrativa e a Organização Acadêmica.
As notas \textbf{brutas} das questões discursivas em cada uma das partes, Formação Geral/Componente Específico, foram calculadas a partir da média simples das notas de cada uma das questões correspondentes.
A primeira questão discursiva da Formação Geral foi desconsiderada e, portanto, a nota \textbf{bruta} referentes às questões discursivas de Formação Geral foi igual à nota da questão mantida, a segunda.
No Componente Específico, a nota \textbf{bruta} correspondeu à média das notas nas três questões.
As notas \textbf{brutas} das questões objetivas de cada uma das partes, Formação Geral/Componente Específico, foram calculadas a partir dos acertos dos estudantes considerando -se uma regra de três para a conversão para a escala de 0 a 100, descontando -se aquelas questões que foram desconsideradas pela comissão ou que não passaram pelo teste do índice de discriminação ponto-bisserial ( ver seção 1.4.6 para a descrição do índice ).
As questões com índices de discriminação inferior a 0,2 foram excluídas do cômputo da nota do estudante.
As notas \textbf{brutas} de Formação Geral e no Componente Específico para cada um dos estudantes foram calculadas como médias ponderadas das suas notas nas questões objetivas e discursivas, ponderação diferente para cada parte : C\_FG = C\_FG\textasciicircum{}OBJ × 0,6 + C\_FG\textasciicircum{}DIS × 0,4 C\_CE = C\_CE\textasciicircum{}OBJ × 0,85 + C\_CE\textasciicircum{}DIS × 0,15 onde C\_FG é a nota \textbf{bruta} da Formação Geral, C\_FG\textasciicircum{}OBJ é a nota \textbf{bruta} nas questões objetivas da Formação Geral, C\_FG\textasciicircum{}DIS é a nota \textbf{bruta} nas questões discursivas da Formação Geral, C\_CE é a nota \textbf{bruta} do Componente Específico, C\_CE\textasciicircum{}OBJ é a nota \textbf{bruta} nas questões objetivas do Componente Específio, e C\_CE\textasciicircum{}DIS é a nota \textbf{bruta} nas questões discursivas do Componente Específico.
A nota \textbf{bruta} da prova, C\_GER, por sua vez, é também uma média ponderada da nota de cada uma das partes, Formação Geral/Componente Específico : C\_GER = C\_FG × 0,25 + C\_CE × 0,75. 3.1 Estatísticas Básicas da Prova 3.1.1 Estatísticas Básicas Gerais A Tabela 3.1 apresenta as estatísticas básicas da prova por grupo de estudantes.
O total de estudantes inscritos para a prova foi de 604.
Destes, 71,2 \% compareceram, sendo que a presença foi maior entre concluintes ( 86,2 \% ) do que entre ingressantes ( 64,0 \% ).
A média geral da prova foi 46,59, sendo que os ingressantes obtiveram média mais baixa, 44,74, que os concluintes, cuja média foi 49,44.
O desvio-padrão geral foi de 16,6, sendo o do grupo de ingressantes menor ( 14,7 ), indicando que, em contraposição ao grupo dos concluintes ( 18,8 ), aqueles possuíram uma distribuição mais homogênea nas notas.
A nota máxima foi 80,6, obtida por pelo menos um concluinte, ao passo que a maior nota obtida por um ingressante foi 80,2.
Os coeficientes de assimetria para os dois grupos considerados, o de concluintes e o de ingressantes, bem como para a população como um todo foram todos negativos, indicando uma maior concentração à esquerda da distribuição e um maior espalhamento à direita.
Uma comparação entre o desempenho de ingressantes e concluintes pode ser realizada a partir dos dados apresentados no Gráfico 3.1.
Cumpre destacar que as notas dos concluintes, de um modo geral, foram mais elevadas do que as dos ingressantes, cuja distribuição de notas está à esquerda dos primeiros.
A ideia de melhor desempenho é reforçada pelas notas medianas dos dois grupos : 52,1 entre os concluintes e 45,4 entre os ingressantes.
Nas faixas de notas até 50 pontos o percentual de ingressantes é maior que o dos concluintes, situação esta que se inverte para as faixas de notas superiores a 50 pontos.
A exceção é o primeiro intervalo de 0 a 10 que inclui uma quantidade maior de concluintes do que de ingressantes.
Constatou -se que aproximadamente 38,7 \% dos ingressantes e 56,2 \% dos concluintes situam -se nas faixas de notas superiores a 50,0 pontos.
O Gráfico 3.2 apresenta informações referentes ao desempenho geral de ingressantes e concluintes, comparando os resultados considerando a Região, Categoria Administrativa e Organização Acadêmica.
As notas médias mais elevadas entre os concluintes ocorreram nas Regiões Sul ( 53,61 ) e Sudeste ( 50,13 ), ambas acima da média nacional ( 49,44 ).
Em relação aos ingressantes, as médias mais altas foram obtidas na Região Sul ( 53,38 ) nas demais Regiões a média foi inferior à média nacional ( 44,74 ).
Na Região Nordeste foi obtida a média de notas mais baixa entre os concluintes ( 43,24 ) e os ingressantes ( 42,75 ).
Cursos nesta área foram oferecidos em instituições federais, estaduais e particulares.
Levando -se em conta os agrupamentos dos estudantes em categorias administrativas, observa -se que a pontuação média mais elevada entre os concluintes e ingressantes foi encontrada em instituições de origem federal ( respectivamente 51,71 e 47,30 ) – ambas acima da média nacional.
O desvio padrão foi menor nas entidades particulares, tanto para ingressantes ( 1,26 ) quanto para concluintes ( 1,57 ).
Tendo como foco as notas médias dos estudantes agrupados por Organização Acadêmica, notamos que só Universidades apresentaram as maiores médias entre os concluintes ( 51,78 ) e Faculdades apresentaram as maiores médias entre os ingressantes ( 50,65 ).
Ao passo que as piores notas médias foram observadas nas Faculdades para concluintes ( 47,33 ) e Universidades para ingressantes ( 35,40 ). 3.1.2 Estatísticas Básicas de Formação Geral A Tabela 3.2 apresenta as estatísticas básicas de Formação Geral por grupo de estudantes.
O total de estudantes inscritos para a prova foi de 604.
Destes, 71,2 \% compareceram, sendo que a presença foi maior entre concluintes ( 86,2 \% ) do que entre ingressantes ( 64,0 \% ).
A média geral da parte da prova de Formação Geral foi de 49,26, sendo que os ingressantes obtiveram média mais baixa, 48,95, que os concluintes, cuja média foi 49,74.
O desvio-padrão geral foi de 23,4, sendo o do grupo de ingressantes menor ( 22,7 ), do que o do grupo dos concluintes ( 24,5 ), aqueles possuíram uma distribuição mais homogênea nas notas.
A nota máxima foi 91,4, obtida por estudantes de ambos os grupos.
Para este componente os coeficientes de assimetria são negativos para o total e os concluintes.
O Gráfico 3.3 permite a comparação entre concluintes e ingressantes relativamente ao desempenho na parte da prova que aborda a Formação Geral.
Nesta parte, concluintes e ingressantes apresentaram um desempenho similar.
A mediana dos concluintes foi de 51,2 enquanto a dos ingressantes foi de 48,8.
Dos ingressantes, aproximadamente 49,0 \% alcançaram nota superior a 50 ; dentre os concluintes, o percentual dos estudantes nesta faixa de notas foi muito semelhante e igual a 52,7 \%.
No Gráfico 3.4 são apresentadas as informações referentes ao desempenho dos ingressantes e dos concluintes na parte de Formação Geral, em diferentes agregações como Região, Categoria Administrativa e Organização Acadêmica.
Considera -se as notas médias dos estudantes segundo cada Região.
As médias variaram, entre os concluintes, de 56,36 na Região Sul a 45,58 na Região Nordeste.
Entre os ingressantes, os valores vão de 42,01 na Região Norte a 57,62 na Região Sul.
A maior dispersão de notas dos estudantes de uma dada Região foram encontrada no Sul para os concluintes e para os ingressantes.
Na Região Sudeste, os concluintes apresentaram médias maiores do que os ingressantes, a situação inversa acontecendo no Nordeste e Sul.
As Regiões Norte e Centro-Oeste não apresentaram concluintes.
Tendo como foco as categorias administrativas, destacou -se que as instituições federais, estaduais e particulares estavam representadas.
Sendo que as notas médias das federais foram superiores para ingressantes ( 55,13 ) e para concluintes ( 55,30 ).
Cabe destacar que entre as instituições estaduais não houve concluintes.
Considerando -se o tipo de Organização Acadêmica, as Universidades apresentaram o maior valor para concluintes ( 52,59 ).
E as Faculdades apresentaram o maior valor para ingressantes ( 52,52 ), sendo este inclusive superior à nota dos seus concluintes.
Os Centros Federais de Educação Tecnológica e os Institutos Federais de Educação, Ciência e Tecnologia não estavam representados nesta área. 3.1.3 Estatísticas Básicas do Componente Específico A Tabela 3.3 apresenta as estatísticas básicas referentes ao Componente Específico da prova, tendo sua média geral, de 45,75, inferior ao desempenho da prova de COMUNICAÇÃO SOCIAL – EDITORAÇÃO como um todo, cuja média foi 46,59, como apresentado na Tabela 3.1, que por sua vez foi inferior ao desempenho na Formação Geral, que teve média igual a 49,26, conforme mostrado na Tabela 3.2.
Os concluintes apresentaram um desempenho médio ( 49,40 ) superior ao dos ingressantes ( 43,39 ), havendo ainda indícios de maior homogeneidade entre as notas médias dos ingressantes ( desvio-padrão de 15,8 ) do que entre as dos concluintes ( 19,3 ), como sugeriu o desvio padrão, menor no grupo dos ingressantes do que no dos concluintes.
A nota máxima obtida por concluintes foi 84,5 e por ingressantes 84,7.
As distribuições de notas para cada grupo se apresentaram com uma inclinação para notas maiores, como evidenciados pelos coeficientes de assimetria negativos.
A situação de assimetria é confirmada pelo Gráfico 3.5, que proporciona uma comparação do desempenho entre concluintes e ingressantes em relação ao Componente Específico da prova.
Dentre as três comparações já feitas ( confrontar com Gráficos 3.1 e 3.3 ) observou -se que, neste componente, a diferença de desempenho entre concluintes e ingressantes foi positivamente a favor dos concluintes.
Claramente a distribuição das notas médias dos ingressantes, à esquerda daquela dos concluintes, evidenciou o percentual mais elevados de notas médias nas faixas mais baixas para este grupo, com exceção do primeiro, que como já comentado, incluiu uma maior proporção de zeros entre os concluintes.
As medianas das distribuições das notas dos ingressantes e dos concluintes foram 44,4 e 50,5, respectivamente, enfatizando o melhor desempenho dos concluintes neste componente da prova.
Assim como os Gráficos 3.2 e 3.4, o Gráfico 3.6 apresenta uma comparação dos resultados em relação às Regiões, Categoria Administrativa e Organização Acadêmica, agora levando em conta o desempenho dos ingressantes e dos concluintes no Componente Específico da prova.
Tendo como foco as notas médias dos estudantes, em determinada Região, observa -se que, entre os concluintes, as mais elevadas foram encontradas nas Regiões Sudeste ( 50,32 ) e Sul ( 52,75 ), valores acima da média nacional para esta parte ( 49,40 ).
Entre os ingressantes, a maior nota ocorreu na Região Sul ( 52,02 ), única acima da média nacional ( 43,38 ).
A Região Sudeste apresenta os contingentes mais homogêneos com respeito ao Componente Específico, tanto entre os ingressantes ( 1,33 ) como entre os concluintes ( 1,81 ).
Novamente, para a análise segundo a Categoria Administrativa, as entidades federais, estaduais e particulares foram representadas.
Observa -se que as médias mais elevadas entre os concluintes foram encontradas em instituições de origem federal ( 50,56 ).
Já entre os ingressantes nas estaduais ( 44,85 ).
Os estudantes das instituições particulares apresentaram uma maior homogeneidade, com desvio padrão das notas maiores para os concluintes do que para os ingressantes.
Quanto à Organização Acadêmica, o comentário é o mesmo já feito para a nota \textbf{bruta} e referente à Formação Geral : somente Universidades, Centros Universitários e Faculdades foram representados neste curso.
Considerando -se o tipo de Organização Acadêmica, entre os concluintes, Centros Universitários apresentaram os maiores valores, seguidos das Universidades e das Faculdades.
E entre os ingressantes, Faculdades apresentaram os maiores valores, seguidas das Universidades e dos Centros Universitários. 3.2 Análise das Questões Objetivas 3.2.1 Formação Geral A Tabela 3.4 apresenta as estatísticas básicas relativas às oito questões objetivas da Formação Geral dos estudantes.
Adiantamos que a primeira questão foi desconsiderada pelo critério da correlação ponto-bisserial.
Como pode ser observado na Tabela 3.4, os estudantes concluintes obtiveram desempenho inferior, com nota média de 56,94, em relação aos estudantes ingressantes cuja nota média foi 57,65.
A variabilidade entre os concluintes foi maior do que entre os ingressantes, como indicaram os desvios-padrão de 20,4, para o grupo de ingressantes e 23,7 para o grupo dos concluintes.
A amplitude total das notas foi a mesma para ambos os grupos : nota mínima de 0,0 e máxima de 85,7.
A mediana foi igual para ambos os grupos : 57,1.
A assimetria, medida pelo coeficiente de assimetria ( skewness ) foi um pouco maior em módulo para os concluintes do que para os ingressantes, ainda que ambos fossem negativos indicando uma inclinação para a direita.
A Tabela 3.5 apresenta a classificação das questões objetivas da Formação Geral, segundo o índice de facilidade ( questões em vermelho foram desconsideradas pelo critério ponto-bisserial ).
Como as questões foram comuns a todas as áreas, a tabela apresenta a classificação considerando -se os concluintes de todas as áreas e aqueles do curso de COMUNICAÇÃO SOCIAL – EDITORAÇÃO.
De acordo com os índices obtidos, as questões objetivas da prova foram assim avaliadas para a área : das oito questões, três tiveram entre 60 \%, inclusive, e 85 \% de acertos e foram classificadas como fáceis ( somente uma para o total de concluintes ) ; uma foi considerada mediana, alcançando entre 40 \%, inclusive, e 60 \% de acertos ( três para o total de concluintes ) ; duas obtiveram entre 15 \%, inclusive, e 40 \% de acertos, sendo consideradas difíceis ( três para o total de concluintes ).
Por fim, apenas uma foi classificada como Muito difícil tendo obtido percentual menor do que 15 \% de acertos ( uma para o total de concluintes ), coincidentemente a questão eliminada pelo critério ponto-bisserial.
Apenas uma questão foi classificada na categoria Muito fácil.
A Tabela 3.6, descreve os resultados das análises das questões objetivas relativas à Formação Geral, segundo o poder de discriminação, utilizando para tal o índice de discriminação ponto-bisserial.
Como as questões foram comuns a todas as áreas, a tabela apresenta a classificação considerando -se os concluintes de todas as áreas e aqueles do curso de COMUNICAÇÃO SOCIAL – EDITORAÇÃO.
A desconsideração da questão segundo este critério considerou o universo de todos os concluintes.
Três das oito questões apresentaram índices iguais ou acima de 0,40 ( nenhuma para a população de concluintes como um todo ), sendo, portanto, classificadas como muito boas nesta característica, para este grupo de estudantes, duas foram classificadas como boas com índice entre 0,30, inclusive, e 0,40 ( três para os concluintes como um todo ) e duas foram classificadas como médias, com índice entre 0,20, inclusive, e 0,30 ( quatro para os concluintes como um todo ).
Apenas uma das questões foi classificada como fraca para os concluintes de modo geral e também para os concluintes de COMUNICAÇÃO SOCIAL – EDITORAÇÃO, sendo desconsiderada pelo critério. de acertos ( três para o total de concluintes ) ; duas obtiveram entre 15 \%, inclusive, e 40 \% de acertos, sendo consideradas difíceis ( três para o total de concluintes ).
Por fim, apenas uma foi classificada como Muito difícil tento obtido percentual menor do que 15 \% de acertos ( uma para o total de concluintes ), coincidentemente a questão eliminada pelo critério ponto-bisserial.
Apenas uma questão foi classificada na categoria Muito fácil.
A Tabela 3.6, descreve os resultados das análises das questões objetivas relativas à Formação Geral, segundo o poder de discriminação, utilizando para tal o índice de discriminação ponto-bisserial.
Como as questões foram comuns a todas as áreas, a tabela apresenta a classificação considerando -se os concluintes de todas as áreas e aqueles do curso de COMUNICAÇÃO SOCIAL – EDITORAÇÃO.
A desconsideração da questão segundo este critério considerou o universo de todos os concluintes.
Três das oito questões apresentaram índices iguais ou acima de 0,40 ( nenhuma para a população de concluintes como um todo ), sendo, portanto, classificadas como muito boas nesta característica, para este grupo de estudantes, duas foram classificadas como boas com índice entre 0,30, inclusive, e 0,40 ( três para os concluintes como um todo ) e duas foram classificadas como médias, com índice entre 0,20, inclusive, e 0,30 ( quatro para os concluintes como um todo ).
Apenas uma das questões foi classificada como fraca para os concluintes de modo geral e também para os concluintes de COMUNICAÇÃO SOCIAL – EDITORAÇÃO, sendo desconsiderada pelo critério.
Na sequência, no Quadro 3.1, são especificadas, por questão, as habilidades aferidas pelas questões de Formação Geral da prova da COMUNICAÇÃO SOCIAL – EDITORAÇÃO, bem como os resultados da análise relativa aos índices de facilidade e de discriminação de cada questão.
O índice de facilidade para a área de COMUNICAÇÃO SOCIAL – EDITORAÇÃO teve uma variação considerável, oscilando entre 0,10 até 0,87 e o de discriminação ficou entre 0,10 e 0,54.
A questão 7 aferiu as habilidades de “ Analisar e criticar informações ”, “ Extrair conclusões por indução e/ou dedução ”, “ Estabelecer relações, comparações e contrastes em diferentes situações ”, “ Detectar contradições ”, “ Fazer escolhas \textbf{valorativas} avaliando consequências ”, “ Questionar a realidade ”, e “ Argumentar coerentemente ” foi a que apresentou, para a população de concluintes de COMUNICAÇÃO SOCIAL – EDITORAÇÃO, maior poder discriminatório com índice de 0,54.
A questão mais fácil, com 87 \% de acertos, foi a de número 8, que aferiu também, além das já listadas, a habilidade de “ Ler e interpretar textos ”.
De outra parte, a questão com menor índice de facilidade e de discriminação deste conjunto foi a questão 1, que foi desconsiderada, e que deveria aferir as habilidades de “ Analisar informações ” ; “ Extrair conclusões ” e “ Propor soluções para situações-problema ”.
O Gráfico 3.7 mostra, a título de ilustração, o comportamento das respostas à questão 7 de formação Geral para os estudantes concluintes da área de COMUNICAÇÃO SOCIAL – EDITORAÇÃO como função da nota geral na prova4.
Trata -se de uma questão considerada fácil para os concluintes de COMUNICAÇÃO SOCIAL – EDITORAÇÃO e de dificuldade média para o conjunto de concluintes de todas as áreas.
O poder de discriminação foi muito bom para os estudantes da área e bom para o conjunto de todos os concluintes.
Neste Gráfico, cada uma das cinco curvas representou o percentual de respostas em determinada alternativa da questão, em função do escore total dos estudantes concluintes da área de COMUNICAÇÃO SOCIAL – EDITORAÇÃO.
A curva azul marinho corresponde à alternativa A, a correta para esta questão.
Assim, observa -se que entre os estudantes com notas mais baixas – desempenho geral mais fraco na prova como um todo –, houve a escolha de alternativas incorretas, especialmente a alternativa C – curva em vermelho seguida da alternativa E – curva em roxo.
A proporção de estudantes que selecionaram a resposta correta A, chegou ao acerto de quase 100 \% dos estudantes com notas entre 62,5 e 75,0, tal percentual declinou para estudantes com nota pelo menos 75,0.
Esta análise permitiu verificar como a questão discriminou entre os grupos de desempenho.
Os Gráficos relativos às demais questões de Formação Geral constam do Anexo I. 3.2.2 Componente Específico A Tabela 3.7 apresenta as estatísticas básicas em relação às questões objetivas do Componente Específico.
Como pode ser observado, os estudantes concluintes obtiveram desempenho superior, com média de 51,91 pontos, em relação aos estudantes ingressantes com uma média de 45,18 pontos.
Os desvios-padrão indicaram que a variabilidade entre os concluintes foi maior ( 19,6 ) do que entre os ingressantes ( 16,1 ).
Os coeficientes de assimetria para os dois grupos foram negativos indicando uma inclinação para a direita.
A assimetria foi, porém, maior entre os concluintes.
A análise das questões objetivas do Componente Específico em COMUNICAÇÃO SOCIAL – EDITORAÇÃO, em relação ao grau de facilidade da prova ( ver Tabela 3.8 ), mostrou que, das 27 questões, 10 encontraram -se na categoria difícil ( entre 15 \% e 40 \%, não incluso, de acerto ), 11 foram de nível médio ( entre 40 \% e 60 \%, não incluso, de acertos ), 4 estavam na categoria fácil ( entre 60 \% e 85 \%, não incluso, de acertos ) e nenhuma teve a classificação muito fácil ( a partir de 85 \%, não incluso, de acertos ).
Nenhuma questão teve percentual de acertos menor do que 15 \% ( classificação muito difícil ).
Duas questões foram desconsideradas pela Comissão Assessora da Área de COMUNICAÇÃO SOCIAL – EDITORAÇÃO.
Das 27 questões objetivas relativas ao Componente Específico ( ver Tabela 3.9 ), sete obtiveram o índice de discriminação muito bom, oito obtiveram o índice de discriminação bom e cinco obtiveram nível de discriminação médio.
O índice fraco que eliminaria a questão, foi alcançado por cinco questões e duas questões foram canceladas pela comissão.
Na sequência, no Quadro 3.2 estão especificadas, por questão, as habilidades aferidas pelas questões do Componente Específico da prova, bem como os resultados da análise relativa aos índices de facilidade e de discriminação de cada questão.
O Gráfico 3.8 mostra, a título de ilustração, o comportamento da questão 11 do Componente Específico para os estudantes concluintes da área de COMUNICAÇÃO SOCIAL – EDITORAÇÃO como função da nota geral na prova.
Trata -se de uma questão considerada fácil para os concluintes de COMUNICAÇÃO SOCIAL – EDITORAÇÃO.
O poder de discriminação foi bom para os estudantes da área.
Neste Gráfico, cada uma das cinco curvas, representa o percentual de respostas em determinada alternativa da questão, em função da nota total dos estudantes concluintes de COMUNICAÇÃO SOCIAL – EDITORAÇÃO.
A curva verde claro correspondeu à alternativa B, a correta para esta questão.
Assim, observou -se que entre os estudantes com escores mais baixos – desempenho geral mais fraco na prova como um todo –, a tendência foi a escolha de uma alternativa incorreta, a alternativa A – curva azul marinho, seguida da alternativa E – curva roxa.
A proporção de estudantes que selecionaram a resposta correta aumentou com o desempenho deste grupo de estudantes chegando ao acerto de quase 90 \% dos estudantes com notas pelo menos 75,0.
Esta análise permitiu verificar como a questão discriminou entre os grupos de desempenho. roxa.
A proporção de estudantes que selecionaram a resposta correta aumentou com o desempenho deste grupo de estudantes chegando ao acerto de quase 90 \% dos estudantes com notas pelo menos 75,0.
Esta análise permitiu verificar como a questão discriminou entre os grupos de desempenho.
Os Gráficos relativos às demais questões do Componente Específico constam do Anexo I. 3.3 Análise das Questões Discursivas 3.3.1 Formação Geral A Tabela 3.10 apresenta as estatísticas básicas da prova por grupo de estudantes.
A análise dos resultados de desempenho dos estudantes nas questões discursivas que são apresentados na tabela 3.10 e no gráfico 3.9, mostrou que as notas foram mais baixas no conjunto dessas questões que no das questões objetivas.
O total de estudantes inscritos para a prova foi de 604.
Destes, 71,2 \% compareceram, sendo que a presença foi maior entre concluintes ( 86,2 \% ) do que entre ingressantes ( 64,0 \% ).
A média geral da prova foi 37,20, sendo que os ingressantes obtiveram média mais baixa, 36,02, que os concluintes, cuja média foi 39,04.
O desvio-padrão geral foi de 38,8, sendo o do grupo de ingressantes menor ( 37,9 ), em relação ao grupo dos concluintes ( 40,2 ), aqueles possuíram uma distribuição mais homogênea nas notas.
A nota máxima foi 100,0, atingida por concluintes e ingressantes.
Os coeficientes de assimetria para os dois grupos considerados, o de concluintes e o de ingressantes, bem como para a população como um todo são todos positivos, indicando uma maior concentração à esquerda da distribuição.
Para todos os grupos, um pouco menos da metade dos estudantes presentes obtiveram nota zero nas questões discursivas da Formação Geral.
Das 206 provas com nota zero, 21 ( 10,2 \% ) continham um protesto, 45 ( 21,8 \% ) estavam em branco e 141 ( 68,4 \% ) foram respondidas, porém de forma inadequada.
Entre os concluintes, a maioria ( 64,2 \% ) dentre aqueles com nota zero, apresentaram as provas respondidas de forma inadequada.
Entre os ingressantes, 71,2 \% das notas nulas decorreram de correção.
Uma comparação entre o desempenho de ingressantes e concluintes pode ser realizada a partir dos dados apresentados no Gráfico 3.9.
As notas dos concluintes, de um modo geral, foram ligeiramente mais elevadas do que as dos ingressantes, cuja distribuição de notas está à esquerda dos primeiros.
A proporção de notas nulas foi menor entre os concluintes ( 46,2 \% ) do que entre os ingressantes ( 46,7 \% ).
A mediana foi maior para os concluintes ( 42,5 ) do que para os ingressantes ( 25,0 ).
Na sequência, discorrer -se-á, especificamente, sobre uma das questões discursivas de Formação Geral, a número 10 já que a número 9 foi desconsiderada.
Em especial, serão analisados os conteúdos englobados na questão e os resultados obtidos pelos estudantes.
Ao final, serão apresentados os comentários da banca corretora a respeito do observado nas respostas dos estudantes, suas impressões e conclusões.
Cumpre esclarecer que, tendo em vista que as questões discursivas de Formação Geral foram padronizadas, ou seja, constam de todas as provas, os comentários da banca foram os mesmos para todas as áreas acadêmicas, eis que direcionados a todos os estudantes que participaram do ENADE/2009.
A seguir, serão analisados os desempenhos de ingressantes e concluintes da área de COMUNICAÇÃO SOCIAL – EDITORAÇÃO na questão discursiva de Formação Geral, não anulada, do ENADE/2009, comparando os resultados obtidos com a habilidade exigida na questão. 3.3.1.1 Análise da Questão Discursiva 9 de Formação Geral DESCONSIDERADA 3.3.1.2 Análise da Questão Discursiva 10 de Formação Geral A Questão Discursiva nº. 10 buscou avaliar as habilidades : • Ler e interpretar textos ; • Analisar e criticar informações ; • Extrair conclusões por indução e/ou dedução ; • Estabelecer relações, comparações e contrastes em diferentes situações ; • Fazer escolhas \textbf{valorativas} avaliando consequências ; • Questionar a realidade ; • Argumentar coerentemente.
O Quadro 3.3 apresenta as habilidades/competências, temas predominantes e estatísticas básicas da questão 10 de Formação Geral por grupo de estudantes.
Cumpre notar que como a questão 9 foi desconsiderada, a nota referente às questões discursiva da Formação Geral é idêntica à nota da questão 10.
A média geral desta questão foi 37,20, sendo que os ingressantes obtiveram média mais baixa, 36,02, que os concluintes, cuja média foi 39,04.
O desvio-padrão geral foi de 38,8, sendo o do grupo de ingressantes menor ( 37,9 ), em relação ao grupo dos concluintes ( 40,2 ), aqueles possuíram uma distribuição mais homogênea nas notas.
A nota máxima foi 100,0, atingida por concluintes e ingressantes.
Os coeficientes de assimetria para os dois grupos considerados, o de concluintes e o de ingressantes, bem como para a população como um todo foram positivos, indicando uma maior concentração à esquerda da distribuição.
Para todos os grupos, um pouco menos da metade dos estudantes presentes tiveram nota zero nesta questão discursiva de Formação Geral.
Das 206 provas com nota zero, 21 ( 10,2 \% ) continham um protesto, 45 ( 21,8 \% ) estavam em branco e 141 ( 68,4 \% ) foram respondidas, porém de forma inadequada.
Entre os concluintes, a maioria ( 64,2 \% ) dentre aqueles com nota zero estavam com as provas respondidas de forma inadequada.
Entre os ingressantes, 71,2 \% das notas nulas( zero ) decorreram do desempenho do próprio estudante verificadas na correção.
Uma comparação entre o desempenho de ingressantes e concluintes para esta questão foi exatamente igual à já realizada para a nota discursiva de Formação Geral a partir dos dados apresentados no Gráfico 3.9. 3.3.1.3 Comentários sobre as respostas à Questão Discursiva 10 Apresentação A questão 10 de Formação Geral apresentou um questionamento capaz de analisar a habilidade do estudante em reconhecer o que é um direito social, optar por um direito social vigente que possa diminuir as desigualdades de renda familiar no Brasil e discutir os argumentos em favor do direito social escolhido como suporte à resposta.
Para análise das respostas apresentadas pelos estudantes, e para efeito de pontuação, foram considerados os seguintes argumentos : • Acesso à educação pública, gratuita e de qualidade, garantindo ao cidadão melhores condições de vida ; • Assistência à saúde, em seu contexto mais amplo, garantindo os cuidados essenciais de saúde, baseados em técnicas apropriadas ao alcance de todas as pessoas, contando com a participação da população ; • Permanência do estudante na escola, em todos os níveis escolares ( da educação infantil à educação superior, possibilitando a qualificação profissional ), o que garante melhores condições de trabalho e remuneração.
Considerou -se a qualificação e o desempenho docente, desenvolvendo nos estudantes melhores expectativas de aprendizagem e interesse pelo conteúdo da escola.
Os professores desenvolvendo metodologias adequadas e conhecimento pleno quanto a questões de aprendizagem ; • Condições dignas de trabalho, com remuneração que garanta qualidade de vida.
Foram consideradas as condições de trabalho, a preparação profissional para o trabalho e os cursos de aperfeiçoamento, oportunizando momentos de aprendizagem e desenvolvimento, a partir de situações desencadeadoras ; • Ser proprietário do imóvel em que reside, reduzindo os gastos com aluguel, garantindo melhores condições de vida, equilíbrio financeiro e construção do bem-estar familiar.
Desempenho dos estudantes na questão É importante salientar que havia respostas que desenvolveram adequadamente o questionamento apresentado, com fundamentação baseada em argumentos coerentes, contudo, em um número considerável de respostas, havia apenas um argumento e não dois, conforme fora solicitado.
Muitos estudantes informaram sobre problemas da educação, de um modo geral.
Algumas observações mencionadas pelos corretores durante o processo de correção foram : • há preocupação com uma educação de qualidade, com sugestões de cursos profissionalizantes e métodos preparatórios de aulas. • a educação é revelada como de baixa qualidade, exaltando -se treinamentos e cursos de atualização para a classe docente. • a saúde como essência para a vida.
O SUS é colocado como parte integrante de direito social, garantindo às pessoas fácil acesso. • oportunidade de emprego como forma de ascensão social, direito, cidadania. • os políticos brasileiros como causadores do despreparo educacional.
As respostas analisadas apresentaram norma culta, coesão e coerência razoáveis, que garantiram a compreensão da mensagem desenvolvida pelos estudantes.
Os erros ortográficos foram poucos e não comprometeram a discussão sobre o assunto e sobre os pontos de vista abordados.
Quanto à coerência, houve um desvio em relação ao objetivo da questão, já que alguns estudantes, através dos seus argumentos retrataram reivindicações, reclamações e situações de problemas enfrentados pelos brasileiros, citando golpes políticos.
Porém, mesmo quando se tratava desta avaliação, os estudantes foram autênticos através das explicações.
Houve erros quanto à concordância verbal e nominal.
Houve, também, dificuldade em utilizar os conectores corretamente.
Alguns estudantes, por falta de ideias, apresentaram em suas respostas redundância e prolixidade.
O desempenho foi considerado satisfatório, pois a maioria dos estudantes soube responder ao questionamento.
Houve domínio do conteúdo, pois as ideias, os argumentos e a \textbf{originalidade} corresponderam ao esperado. 3.3.2 Componente Específico Nas questões discursivas do Componente Específico – questões de números 38, 39 e 40 –, a diferença de desempenho entre os grupos não foi muito acentuada, inclusive menor do que para as questões discursivas de Formação Geral.
A Tabela 3.11 contém as estatísticas da média destas questões.
Observa -se que no grupo dos concluintes, a nota média alcançada foi 35,49, enquanto que no grupo dos ingressantes foi 33,57, ambas menores do que as obtidas nas questões discursivas de Formação Geral, que foram 36,02 para os ingressantes e 39,04 para os concluintes ( Tabela 3.10 ).
As médias foram também menores do que as obtidas para as questões objetivas desse componente ( Tabela 3.7 ) : respectivamente 45,2 e 51,9 para ingressantes e concluintes.
A análise de cada uma destas questões será feita nas subseções seguintes.
Os valores positivos para a medida de assimetria indicaram uma distribuição com cauda para a esquerda.
Dos 430 presentes, 29,1 \% obtiveram nota zero.
A tabela 3.12 apresenta um detalhamento da situação destes estudantes.
A proporção dos que não fizeram nenhuma das 3 questões foi de 39,0 \%.
Esta proporção foi menor entre os ingressantes ( 38,5 \% ) do que entre concluintes ( 40,0 \% ).
A proporção de todas as questões em branco correspondeu quase 40 \% das provas.
A proporção de todas as questões respondidas correspondeu quase 20 \% das provas.
Menos de 1 \% no total das provas correspondem a 3 questões com protesto.
O Gráfico 3.10 apresenta a distribuição das notas nas questões discursivas do Componente Específico por grupo de estudantes.
A moda se situou no primeiro intervalo, aquele com valores entre 0 e 10 incluindo os extremos do intervalo.
A impressão de assimetria informada pelos coeficientes de assimetria positivos ( Tabela 3.11 ) foi reforçada no gráfico pela concentração no primeiro intervalo.
O Gráfico 3.10 apresenta a distribuição das notas nas questões discursivas do Componente Específico por grupo de estudantes.
A moda se situou no primeiro intervalo, aquele com valores entre 0 e 10 incluindo os extremos do intervalo.
A impressão de assimetria informada pelos coeficientes de assimetria positivos ( Tabela 3.11 ) foi reforçada no gráfico pela concentração no primeiro intervalo. 3.3.2.1 Análise da Questão Discursiva 38 do Componente Específico A questão 38, cujas habilidades aferidas encontram -se descritas no Quadro 3.4, que ainda contém outras informações, seguido da Tabela 3.13 e do Gráfico 3.11, foi aquela em que os grupos de estudantes observados obtiveram um desempenho intermediário entre as questões discursivas do Componente Específico.
A nota média do grupo dos ingressantes ( 34,5 ) foi menor do que a dos concluintes ( 34,8 ).
O percentual de notas zero obtidas nos dois grupos foi de 32,3 \%, maior entre os ingressantes ( 34,1 \% ) do que entre os concluintes ( 29,6 \% ).
Mais da metade dos concluintes e dos ingressantes que tiraram zero deixaram a questão em branco ( ver Tabela 3.13 ).
Uma minoria ( 4,3 \% ) utilizou a questão para protestar.
Apesar do baixo desempenho, em ambos os grupos a nota máxima de 100,0 pontos foi alcançada.
Os coeficientes de assimetria positivos indicam uma inclinação para a esquerda e cauda para a direita.
O Gráfico 3.11 apresenta a distribuição das notas na questão 38 discursiva do Componente Específico por grupo de estudantes.
A moda se situou no primeiro intervalo, aquele com valores entre 0 e 10 incluindo os extremos do intervalo.
A impressão de assimetria informada pelos coeficientes de assimetria positivos ( Quadro 3.4 ) foi reforçada no gráfico pela concentração no primeiro intervalo.
A distribuição apresentou outros máximos locais no terceiro e no oitavo intervalo. 3.3.2.2 Comentários sobre as respostas à Questão Discursiva 38 O objetivo central da questão foi baseado na percepção dos estudantes sobre as possibilidades de aproveitamento das campanhas massivas dos grandes veículos de comunicação ( em especial cinema e TV ) para alavancagem de vendas no mercado editorial.
As soluções sugeridas na chave de resposta estavam baseadas no que pode ser observado como prática do mercado editorial, especialmente nos produtos direcionados ao consumo pela faixa etária dos próprios estudantes.
Raríssimas foram as respostas com um bom conteúdo e que obtiveram melhores notas.
A maioria das propostas apenas esboçou soluções inviáveis ( por envolverem despesas elevadas ) ou se limitou a estabelecer, de forma sintética e em uma única linha, a citação de algum evento recente ocorrido no mercado, como o relançamento da série de livros Crepúsculo, Lua Nova e Eclipse de Stephener Meyer, sem qualquer avaliação crítica das ações efetuadas.
Concluindo, podemos afirmar que o enfoque foi de agrado dos estudantes e não foram registrados impropérios ou desabafos, quando da resposta a questão. 3.3.2.3 Análise da Questão Discursiva 39 do Componente Específico A questão 39, cujas habilidades aferidas encontram -se descritas no Quadro 3.5, que ainda contém outras informações, seguido da Tabela 3.14 e do Gráfico 3.12, foi aquela em que os grupos de estudantes observados obtiveram o desempenho intermediário entre as questões discursivas do Componente Específico.
A nota média do grupo dos ingressantes foi de 39,5, e a dos concluintes 36,7.
O percentual de notas zero obtidas nos dois grupos foi de 39,8 \% entre os ingressantes e 39,1 \% entre os concluintes.
Das provas com notas nulas nesta questão, mais da metade foi deixada em branco ( 55,3 \% ) : 67,3 \% entre os ingressantes e 36,4 \% entre os concluintes ( ver Tabela 3.14 ).
Entre aqueles com nota nula, uma minoria ( 4,7 \% ) utilizou a questão para protestar.
Apesar do baixo desempenho, em ambos os grupos a nota máxima de 100,0 pontos foi alcançada.
Os coeficientes de assimetria positivos indicaram uma inclinação para a esquerda e cauda para a direita.
O Gráfico 3.12 apresenta a distribuição das notas na questão 39 discursiva do Componente Específico por grupo de estudantes.
A moda se situou no primeiro intervalo, aquele com valores entre 0 e 10 incluindo os extremos do intervalo.
A impressão de assimetria informada pelos coeficientes de assimetria positivos ( Quadro 3.5 ) foi reforçada no gráfico pela concentração no primeiro intervalo. 3.3.2.4 Comentários sobre as respostas à Questão Discursiva 39 Embora seja considerada uma questão de nível mediano de dificuldade, as respostas apresentadas pelos estudantes estiveram abaixo das expectativas, com a exposição – na maioria dos casos – de poucos conhecimentos acerca do mercado editorial e do tratamento que tais publicações geralmente receberam nas editoras em casos semelhantes ao apresentado.
O objetivo central da questão era a geração de ações que auxiliassem na vinculação entre os títulos de uma coleção com a obra completa de determinado autor, e que influenciassem na sua aquisição como um todo e não de títulos em isolado.
O enunciado era claro e não propiciou qualquer tipo de confusão, ficando o estudante limitado apenas pelo seu conhecimento técnico e sua capacidade de observação dos casos concretos de mercado editorial.
Mais uma vez, o desempenho foi abaixo da média geral.
Grande parte das respostas mal contemplou a repetição dos elementos visuais mais óbvios já citados no enunciado, como projeto gráfico ou fixação de texto.
A referência aos dados mais complexos apontados na chave de resposta – como posfácios, cronologias e análise de enredo – foi rara.
A grande maioria das respostas foi de frases curtas, de menos de três linhas, com pequenas relações de dados sem muito nexo ou reproduzindo parte do enunciado e agregando uma ou outra informação como, por exemplo, “ Faria capas e projetos gráficos idênticos para todos os livros da coleção ”. 3.3.2.5 Análise da Questão Discursiva 40 do Componente Específico A questão 40, cujas habilidades aferidas encontram -se descritas no Quadro 3.6, que ainda contém outras informações, seguido da Tabela 3.15 e do Gráfico 3.13. a questão foi aquela em que os grupos de estudantes estudados observados obtiveram o pior desempenho entre as questões discursivas do Componente Específico.
A nota média do grupo dos ingressantes foi 26,7, e a dos concluintes 35,1.
O percentual de notas nulas nesta questão foi : 54,0 \% entre os ingressantes e 42,6 \% entre os concluintes.
Aqui, a maioria das provas com nota zero decorre de provas com a questão mal formulada : 65,2 \% entre os ingressantes e 58,3 \% entre os concluintes ( ver Tabela 3.15 ).
Uma pequena minoria ( 1,4 \% ), dentre aqueles com nota nula, utilizou a questão para protestar.
Em ambos os grupos, a nota máxima de 100,0 pontos foi alcançada.
Os coeficientes de assimetria positivos indicaram uma inclinação para a esquerda.
O Gráfico 3.13 apresenta a distribuição das notas na questão 40 discursiva do Componente Específico por grupo de estudantes.
A moda se situou no primeiro intervalo, aquele com valores entre 0 e 10 incluindo os extremos do intervalo.
A impressão de assimetria informada pelos coeficientes de assimetria positivos ( Quadro 3.6 ) foi reforçada no gráfico pela concentração no primeiro intervalo. 3.3.2.6 Comentários sobre as respostas à Questão Discursiva 40 O objetivo central da questão era apurar o conhecimento dos estudantes sobre as alterações provocadas pela introdução dos meios digitais de produção e leitura de textos.
O enunciado era focado na questão das diferenças dos hábitos físicos e sociais de leitura nos dois meios, o convencional e o digital, e solicitava do estudante a comparação entre essas duas modalidades.
Essa foi a questão com aproveitamento menos deficiente, porém ainda assim, o desempenho geral foi deficiente na média.
Ainda que a questão abordasse um assunto de vivo interesse e de interação diária dos estudantes, muitos não souberam simplesmente correlacionar suas experiências diárias de leitura para compor o conteúdo das respostas.
A questão do conforto de leitura, mais frequente nas discussões sobre o assunto, foi apontada com certa assiduidade, o que aferiu melhores notas nessa questão.
No entanto, as respostas quase sempre se limitavam a esse aspecto, com textos mais bem elaborados ou de pior qualidade.
As possibilidades de utilização de recursos visuais de maior abrangência nas publicações digitais ( por exemplo : o uso da ferramenta de zoom, que permite uma leitura mais confortável por pessoas com deficiências visuais mais graves ) pouco foram citadas nas respostas, ainda que muitos dos estudantes as utilizassem no seu dia-a-dia.
Aspectos ligados às diferenças nos tratamentos dos projetos gráficos do livro impresso e do livro eletrônico praticamente não ocorreram nas respostas.
O erro mais frequente nesta questão foi a simples listagem de aspectos técnicos da preparação de arquivos para publicações digitais.
O desempenho dos estudantes da área de Editoração não foi bom, em que pese a boa formulação das questões e a relativa simplicidade das chaves de correção, bem como a abordagem de assuntos que perpassam diariamente as vidas dos estudantes em todos os enunciados.
Como o efetivo de estudantes desta habilitação é realmente pequeno, mesmo no conjunto nacional dos cursos, e como uma avaliação mais apurada ficou prejudicada pela baixa ocorrência de provas com respostas relativas a questões desta habilitação e a sua dispersão por uma banca com mais de uma dezena de corretores, o relato e análise do desempenho dos estudantes ficou pouco consubstanciado.

% matched lemmas: bruto, originalidade, valorativo
\end{textsample}
